\documentclass[11pt,a4paper]{article}

\usepackage{amsmath,amssymb}
\usepackage{mathtools}
\usepackage{hyperref}
\usepackage{booktabs}
\usepackage[margin=1.15in]{geometry}
\usepackage{xcolor}

% Shortcuts
\newcommand{\Gstar}{{G^{*}}}
\newcommand{\vpi}{\varpi}
\DeclareMathOperator{\csch}{csch}

% Epistemic margin notes
\newcommand{\etag}[1]{\marginpar{\raggedright\tiny\textsf{#1}}}

% Running header
\usepackage{fancyhdr}
\pagestyle{fancy}
\fancyhf{}
\fancyhead[R]{\small\itshape One Unit of Existence --- Narrative Companion}
\fancyfoot[C]{\thepage}
\renewcommand{\headrulewidth}{0.4pt}

\title{\textsc{One Unit of Existence}\\[6pt]
\large The Life and Death of a Hydrogen Atom\\
as Told by the Constants That Made It}

\author{}
\date{}

\begin{document}
\maketitle
\thispagestyle{empty}

\begin{quote}\itshape
``What is a hydrogen atom?'' is a question that admits an answer at every
level of description. This is the answer at the level where the description
itself becomes the thing described.
\end{quote}

\noindent\textbf{A note to the reader.} This document tells the same story as the formal treatment~\cite{formal}, but without numbered theorems or proofs. The equations that appear here are the same equations proved there. Where a claim's epistemic status matters---where honest accounting demands we say ``this is assumed, not derived''---we say so in the text. For the complete derivation chain, see~\cite{chain}; for the mathematical structure of the discrete--continuous bridge, see~\cite{bridge}; for the ontological constructor unifying these results, see~\cite{constructor}; for the computational engine, see~\cite{engine}.

%=======================================================================
\section*{I.\quad Before the Beginning}
%=======================================================================

Before there is a hydrogen atom, there is a lattice. Not a lattice made of anything, but the lattice that makes anything possible.

The framework begins with a single commitment: reality consists of sites arranged on a cubic grid, $\mathbb{Z}^3$, and each site is in one of exactly three states---negative, void, or positive. Written as numbers: $\{-1,\, 0,\, +1\}$. That is the entire ontology. Two properties per site---state and position---and everything else must emerge.

Each site also carries a continuous vector field $\mathbf{J} \in \mathbb{R}^3$, called the \emph{flux}. The flux encodes tendency: what the site \emph{would do} if conditions for manifestation were met. The ternary state encodes actuality: what the site \emph{is}. This two-layer structure---disposition and actuality, continuous potential and discrete fact---is the architecture from which all physics will be built.

The void state $s = 0$ is not emptiness. It is the cancellation $0 = (-1) + (+1)$: perfect balance, every tendency met by its complement.

Nothing has happened yet. No particles, no fields, no time. Only the lattice and a figure-eight curve that has been waiting since 1694.

%=======================================================================
\section*{II.\quad The Figure Eight}
%=======================================================================

The lemniscate of Bernoulli, $r^2 = \cos 2\theta$, is the simplest closed curve that crosses itself. Its half-perimeter is a number called the lemniscate constant:
\[
\vpi \;=\; 2\int_0^1 \frac{dx}{\sqrt{1 - x^4}} \;=\; 2.62206\ldots
\]
This integral requires only the integer~4. It is the lemniscatic analogue of~$\pi$: where $\pi$ measures how far you must walk around a circle, $\vpi$ measures how far you must walk around a figure eight.

From~$\vpi$, a second constant is born---the \emph{bridge constant}:
\[
\Gstar \;=\; \frac{\Gamma(1/4)^2}{\sqrt{2}\;\Gamma(1/2)^2} \;=\; \frac{2\vpi}{\sqrt{\pi}} \;=\; 2.95868\ldots
\]
The first form contains no~$\pi$: both $\Gamma(1/4)$ and $\Gamma(1/2)$ are definable by integrals over $e^{-t}$ and rational powers of~$t$, with no circular functions. The identity $\Gamma(1/2)^2 = \pi$ is a theorem---a consequence of the Gaussian integral---not an assumption. From $\vpi$ and $\Gstar$, neither of which requires $\pi$ for its computation:
\[
\pi \;=\; \frac{4\,\vpi^2}{\Gstar^2}\,.
\]
This is not circular. Two independently computable transcendental numbers produce a third. The three constants $\vpi$, $\Gstar$, and $\pi$ form a three-way invariant: each expressible through the other two, but only $\pi$ lacks a simple independent definition free of the other two.

The bridge constant is not arbitrary. It is an invariant of the elliptic curve $E: y^2 = x^3 - x$, which has the special property $j = 1728$ and endomorphism ring $\mathbb{Z}[i]$---the Gaussian integers. Among all elliptic curves, this one has the maximal symmetry compatible with a square lattice. (Why this curve and not another? That is an assumption---the first of five selection principles on which the framework depends. The formal treatment~\cite{formal} catalogs and critiques all five.)

The connection to the cubic lattice is through the Watson integral---the lattice Green's function of $\mathbb{Z}^3$---which satisfies $\Gstar = \sqrt{2\pi\,W_3}$. The geometry of the lemniscate and the geometry of the cubic lattice are the same geometry, seen from different angles.

%=======================================================================
\section*{III.\quad Two Numbers from One Equation}
%=======================================================================

The bridge constant now produces an equation:
\[
x^2 \;-\; 16\,\Gstar^2\, x \;+\; 16\,\Gstar^3 \;=\; 0\,.
\]
The coefficient~16 is $|\mathrm{Aut}(E)|^2 = 4^2$, an intrinsic arithmetic invariant of the elliptic curve. (Why this particular polynomial? That is the second selection principle.)

This equation has two positive real roots:
\[
x_+ = 137.036\,171\ldots \qquad\qquad x_- = 3.023\,964\ldots
\]

The larger root is within 1.26 parts per million of the inverse fine-structure constant, $\alpha^{-1} = 137.035\,999\,177(21)$---the number that determines the strength of the electromagnetic force. The integer part of the smaller root is~3: the number of color charges in quantum chromodynamics.

Whether this agreement is coincidence or structure is the central question of the framework. It is labeled a conjecture, not a theorem. But if it is not coincidence, then a single algebraic object---a quadratic over the lemniscate constant---controls both the electromagnetic and strong forces.

%=======================================================================
\section*{IV.\quad Why Three Dimensions}
%=======================================================================

The Watson integral can be computed in any dimension. In $D = 2$, the smaller root gives $\lfloor x_- \rfloor = 8$, not~2. In $D = 4$ through $D = 6$, it gives~2, not~$D$. Only in $D = 3$ does the number of color charges equal the number of spatial dimensions:
\[
\lfloor x_- \rfloor \;=\; 3 \;=\; D\,.
\]
This is a self-referential selection: the dimension determines the lattice geometry, which determines the bridge constant, which determines the gap equation, whose smaller root gives $N_c = D$. Only $D = 3$ closes the loop. Three spatial dimensions are not an axiom of the framework. They are its unique self-consistent solution.

%=======================================================================
\section*{V.\quad The Proton: Made of Motion}
%=======================================================================

The universe cools. At roughly $10^{-5}$ seconds after whatever initiated the expansion, the temperature drops below $\sim 150$~MeV---the QCD confinement scale. Three quarks are confined into a proton.

The confinement is not a cage imposed from outside. It is a topological fact about the $\mathrm{SU}(3)$ color field: stretching the flux tubes between quarks costs more energy than creating new quark-antiquark pairs. The system chooses the bound state.

The number of colors is not arbitrary. In FTD, the 26 neighbors of a lattice site decompose into three sublattices by distance: 6 face-neighbors, 12 edge-neighbors, 8 corner-neighbors. Each sublattice couples to a different number of flux components---one, two, or three---and this determines the gauge symmetry:
\begin{center}
\begin{tabular}{cccc}
\toprule
Sublattice & Count & Flux components & Gauge group \\
\midrule
Face (SC) & 6 & 1 & $\mathrm{U}(1)$ \\
Edge (FCC) & 12 & 2 & $\mathrm{SU}(2)$ \\
Corner (BCC) & 8 & 3 & $\mathrm{SU}(3)$ \\
\bottomrule
\end{tabular}
\end{center}
The Standard Model gauge group $\mathrm{U}(1) \times \mathrm{SU}(2) \times \mathrm{SU}(3)$ is the unique decomposition of the flux field's energy by the Moore neighborhood. It is not postulated. It is counted.

The proton's mass is approximately 938~MeV. Almost none of this comes from the quarks themselves. Over 99\% is the kinetic and binding energy of the strong field, whose confinement scale is set by the QCD beta function coefficient $b_3 = (11 \times 3 - 2 \times 6)/3 = 7$---a framework integer. The proton is, in this precise sense, made of motion, not of matter.

The proton is already a bridge between the discrete and the continuous: a discrete, countable object (baryon number~$= 1$) whose internal structure is a continuous quantum field theory. It will outlive every star.

%=======================================================================
\section*{VI.\quad The Atom Is Born}
%=======================================================================

For the first 380,000 years, the universe is a plasma. Then the temperature drops below $\sim 3000$~K, the photon bath no longer carries enough energy to ionize hydrogen, and nearly every proton captures an electron. The universe becomes transparent.

This is the moment our hydrogen atom is born.

The electron mass in FTD is $m_e = m_P\sqrt{2\pi}\,(16/3)\,\alpha^{11} = 0.5100$~MeV---within 0.20\% of the measured value. (This formula depends on the framework integers $\{3,\,4\}$ and is classified as a selection, not a derivation.) The Higgs quartic coupling $\lambda = 3/23$ emerges from the ternary state decomposition, giving a Higgs boson mass of $m_H = 125.69$~GeV (0.35\% from the measured 125.25~GeV).

The electron does not orbit the proton like a planet. It occupies a standing wave described by the Schr\"odinger equation, with energy levels
\[
E_n \;=\; -\frac{\alpha^2\,m_e c^2}{2n^2}\,, \qquad n = 1,\,2,\,3,\,\ldots
\]
Using the FTD-derived values of $\alpha$ and $m_e$, the ground-state binding energy is $|E_1| = 13.58$~eV, 0.13\% below the measured 13.598~eV---the same relative error as $m_e$ itself, inherited rather than independent.

The entire hydrogen spectrum---every line in the Balmer series, every transition that Bohr, Sommerfeld, and Dirac painstakingly catalogued---follows from a single root of the master quadratic.

%=======================================================================
\section*{VII.\quad The Vacuum Speaks}
%=======================================================================

Our hydrogen atom sits in its ground state. It is silent, invisible, inert. Then a photon arrives.

If the photon carries exactly the right energy---$E_2 - E_1 = 10.2$~eV---the atom absorbs it and the electron jumps to $n = 2$. Not approximately. Exactly. The atom either absorbs the photon or it does not. There is no half-transition. Discreteness is absolute.

But the vacuum is not silent. Even in its ground state, the electromagnetic field fluctuates: virtual photons are constantly created and annihilated. These vacuum fluctuations shift the $2S_{1/2}$ state upward by about 1058~MHz relative to $2P_{1/2}$---the Lamb shift, measured by Willis Lamb in 1947. It scales as $\alpha^5$: five powers of the bridge constant's imprint. Each power represents one more layer of interaction between the discrete and the continuous.

The hyperfine transition---1420.405~MHz, the 21-cm line---scales as $\alpha^4(m_e/m_p)$. The mass ratio introduces the strong force (since $m_p$ is set by QCD), so the hyperfine line is where electromagnetic and strong interactions meet, both encoded in the master quadratic.

Every stage of the atom's existence exhibits the same architecture:

\begin{center}
\begin{tabular}{lll}
\toprule
\textbf{Event} & \textbf{Discrete} & \textbf{Continuous} \\
\midrule
Quark confinement & baryon number & QCD flux tubes \\
Recombination & electron capture & photon field \\
Absorption/emission & energy levels $E_n$ & electromagnetic field \\
Lamb shift & atomic state & vacuum fluctuations \\
Hyperfine transition & spin states & radio-frequency field \\
Stellar fusion & nuclear binding & thermal plasma \\
Proton decay & baryon number $\to 0$ & radiation field \\
\bottomrule
\end{tabular}
\end{center}

At every transition, the bridge between the discrete side and the continuous side is controlled by coupling constants determined by $\vpi$, $\Gstar$, and the integers $\{3,\,4,\,7,\,13\}$.

The bridge is not $\pi$. The bridge is $\Gstar$.

%=======================================================================
\section*{VIII.\quad The Entangled Pair}
%=======================================================================

Two hydrogen atoms, created together in a singlet state, are separated and measured along independently chosen axes. The correlations between their outcomes violate Bell's inequality: no theory in which each atom carries predetermined values can reproduce the data.

How does a deterministic lattice produce nonclassical correlations?

The answer is the Gauss constraint---but not in the way originally expected. The flux field $\mathbf{J}$ has three components, and the constraint $\nabla \cdot \mathbf{J} = \rho$ eliminates one degree of freedom. The physical flux lives in a two-dimensional transverse plane. A local hidden-variable model with sign-function detectors on this 2D plane produces a triangle-wave correlation $E(\theta) = 1 - 2|\theta|/\pi$, giving $S = 2$---exactly the classical Bell bound. No local model can exceed this.

The engine nevertheless produces $S = 2\sqrt{2}$ in simulation. The mechanism is not local. The Gauss constraint is enforced each tick by solving the Poisson equation $\nabla^2\phi = \nabla \cdot \mathbf{J} - s$ globally across the entire lattice. When one particle's state changes, the constraint enforcement instantaneously redistributes flux everywhere---including at the distant entangled partner's location. This is a deterministic, non-local hidden variable theory, structurally analogous to the quantum potential in Bohmian mechanics. The Bell violation is real; its source is the global constraint solver, not a local integral.

(Whether the global Poisson solver represents physical non-locality or is a computational artifact is an open question. A local iterative solver with finite propagation speed would test whether Bell violation persists without instantaneous constraint enforcement.)

%=======================================================================
\section*{IX.\quad Gravity, Stars, and the Arrow of Time}
%=======================================================================

Our hydrogen atom drifts for hundreds of millions of years, part of a diffuse cloud in intergalactic space. Then gravity gathers the cloud. It contracts, heats, and eventually the core temperature reaches $\sim 10^7$~K.

In FTD, the projection from the infinite-dimensional space of smooth vector fields onto the finite lattice is necessarily lossy---no map from an infinite-dimensional space to a finite set can be injective. Information is destroyed at each update cycle. This irreversibility is the arrow of time: entropy increases because the projection discards information that can never be recovered.

Gravity is conjectured to arise from computational bandwidth saturation: where manifested particles cluster densely, the lattice's finite update rate creates a local processing deficit that bends the effective metric. The full Einstein field equations follow via the Deser iterative bootstrap---demanding that the gravitational field's own energy appear as source---and Lovelock's uniqueness theorem in four-dimensional spacetime.

At $\sim 10^7$~K, protons fuse. The tunneling rate through the Coulomb barrier depends exponentially on the Gamow energy, in which $\alpha$ appears again. The value of $\alpha$ from the master quadratic happens to fall in the narrow window compatible with stellar lifetimes of $\sim 10^{10}$ years---long enough for complexity to emerge on their planets.

%=======================================================================
\section*{X.\quad Death and Dissolution}
%=======================================================================

Two protons fuse: one undergoes beta-plus decay, and the result is a deuteron, a positron, and a neutrino. Our hydrogen atom's proton has become part of a deuterium nucleus. The atom, as a distinct entity, is gone.

In more massive stars, the chain continues: helium to carbon, carbon to oxygen, up to iron at the bottom of the valley of stability. Each step is a discrete nuclear transition in a continuous plasma, with rates set by coupling constants from the integers $\{3,\,4,\,7,\,13\}$. The entire periodic table is a consequence of these four numbers.

If the star is massive enough, the iron core collapses. The supernova disperses all elements heavier than helium back into the interstellar medium. Among the debris: new hydrogen atoms. The cycle is ready to begin again.

And if baryon number is not exactly conserved---if the proton eventually decays, in an era when every star has burned out and every black hole has evaporated---then the last discrete object crosses the final bridge into the continuum of radiation. Photons have no rest mass, no charge, no baryon number. They are pure field, pure wave, pure continuity.

The integer becomes the irrational. The lattice point dissolves into the plane.
\[
\sum \;\longrightarrow\; \int
\]

%=======================================================================
\section*{XI.\quad What This Framework Claims---and What It Does Not}
%=======================================================================

The framework contains approximately 40~genuine derivations (from $\Gstar$ and integers alone), 50~parametric insertions (FTD values plugged into standard physics formulas), and 50+ external results adopted without derivation. The hydrogen spectrum is a parametric insertion: the formula $E_n = -\alpha^2 m_e c^2/(2n^2)$ is standard quantum mechanics; FTD provides the inputs.

The five selection principles on which the strongest results depend are:
\begin{enumerate}
\item Why the lemniscatic curve $E: y^2 = x^3 - x$ with $j = 1728$?
\item Why a quadratic polynomial?
\item Why coefficient 16?
\item Why $\varepsilon = e^\pi - \pi - 20$ as the expansion parameter?
\item Why these specific correction coefficients from $\{3,4,7,13\}$?
\end{enumerate}
Each is motivated by symmetry, minimality, or arithmetic structure, but none is uniquely forced. The formal treatment~\cite{formal} states, justifies, and critiques each one.

The framework is falsified if: (1)~precision measurements of $\alpha$ diverge from $x_+ = 137.03617\ldots$ beyond 10~ppm; (2)~a fourth fermion generation is discovered; (3)~no ensemble of local deterministic lattice states can produce Bell correlations $S > 2$; (4)~Lorentz violation with the wrong sign is observed; (5)~the Higgs quartic is measured far from $3/23 = 0.1304$.

A framework that cannot be wrong is not physics.

%=======================================================================
\section*{Coda: The Ratio Remains}
%=======================================================================

The ratio $\ell = \vpi^2/\Gstar^2$ was there before the atom, is there during the atom, and will be there after the atom. The atom is a transient excitation. The constants are not made of atoms.

The hydrogen atom does not ``know'' about the lemniscatic lattice. It does not need to. It simply \emph{is} what the lattice describes, projected into three spatial dimensions and one temporal one, bound by a coupling constant that is a root of a quadratic over~$\Gstar$, radiating at frequencies that are integer multiples of $\alpha^2 m_e c^2$, and destined to dissolve back into the continuum from which it was projected.

\bigskip

\begin{center}\itshape
The bridge is not $\pi$. The bridge is $\Gstar$.\\[3pt]
And the story of one hydrogen atom is the story of every discrete thing\\
that ever existed in a continuous universe.
\end{center}

\medskip

\begin{center}
\textit{Foundational Ternary Dynamics: a framework, not a faith.}
\end{center}

\vfill

\begin{thebibliography}{99}

\bibitem{codata} E.~Tiesinga \textit{et al.}, ``CODATA recommended values of the fundamental physical constants: 2022,'' \textit{Rev.\ Mod.\ Phys.}\ \textbf{96} (2024), 015002.

\bibitem{lamb} W.~E.~Lamb and R.~C.~Retherford, ``Fine structure of the hydrogen atom by a microwave method,'' \textit{Phys.\ Rev.}\ \textbf{72} (1947), 241--243.

\bibitem{superk} Super-Kamiokande Collaboration, ``Search for proton decay via $p \to e^+\pi^0$ and $p \to \mu^+\pi^0$,'' \textit{Phys.\ Rev.\ D}\ \textbf{95} (2017), 012004.

\bibitem{formal} ``One Unit of Existence: The Lifecycle of a Hydrogen Atom from Lattice Geometry to Dissolution,'' companion paper (2026).

\bibitem{chain} ``The Ontic Derivation Chain: Twenty-Six Physical Constants from Two Inputs,'' companion paper (2026).

\bibitem{bridge} ``The Discrete--Continuous Bridge Through the Lemniscatic Lens,'' companion paper (2026).

\bibitem{constructor} ``The Instantiating Formula: An Ontological Constructor for Foundational Ternary Dynamics,'' companion paper (2026).

\bibitem{engine} ``The Lattice Engine: A Computational Implementation of Foundational Ternary Dynamics,'' companion paper (2026).

\end{thebibliography}

\end{document}
